% Adapted from academic-drawing/templates/tikz_template.tex.
% Main ports: south -> north. Feedback remains left of the main column.
\begin{tikzpicture}[
  every node/.style={font=\figurefont\small,align=center,text=BookInk},
  block/.style={book node,font=\figurefont\small,text width=44mm,minimum height=11mm},
  flow/.style={book arrow,rounded corners=1.2mm},
  feedback/.style={flow,dashed,draw=BookAmber}
]
  \node[block,draw=FigureInput,fill=FigureInput!10] (spec)
    at (-45mm,0) {规格与允许变化\\输入、输出、保持条件};
  \node[block,draw=FigureInput,fill=FigureInput!10] (assets)
    at (45mm,0) {已有资产\\仓库版本、接口与依赖};
  \node[block,draw=FigureModel,fill=FigureModel!10] (candidate)
    at (0,-25mm) {生成候选或局部补丁\\保留版本与修改位置};
  \node[block] (test) at (0,-47mm)
    {在固定环境运行\\目标测试与相关回归};
  \node[block] (result) at (0,-69mm)
    {读取实际结果\\返回值、异常、失败断言};
  \node[block,text width=27mm,draw=BookAmber,fill=BookAmber!8]
    (revise) at (-61mm,-69mm) {定位并修订\\计入总预算};
  \node[block,draw=FigureOutput,fill=FigureOutput!10] (freeze)
    at (0,-91mm) {按预定规则选定\\冻结候选与环境};
  \node[block,text width=38mm] (stop) at (60mm,-69mm)
    {停止并报告\\预算耗尽或环境失败};
  \node[block,draw=FigureOutput,fill=FigureOutput!10] (heldout)
    at (0,-118mm) {独立验收\\留出测试与结果记录};
  \coordinate (inputs) at (0,-14mm);
  \draw[draw=BookMuted,line width=.7pt,rounded corners=1.2mm]
    (spec.south) |- (inputs);
  \draw[draw=BookMuted,line width=.7pt,rounded corners=1.2mm]
    (assets.south) |- (inputs);
  \draw[flow] (inputs) -- (candidate.north);
  \draw[flow] (candidate.south) -- (test.north);
  \draw[flow] (test.south) -- (result.north);
  \draw[feedback] (result.west) -- node[above] {未通过} (revise.east);
  \draw[feedback] (revise.north) |- (candidate.west);
  \draw[flow] (result.east) -- (stop.west);
  \draw[flow] (result.south) -- node[right] {通过} (freeze.north);
  \draw[flow] (freeze.south) -- (heldout.north);
  \draw[draw=BookMuted,line width=.8pt,dotted]
    (-78mm,-104mm) -- (80mm,-104mm);
  \node[anchor=east,align=right] at (79mm,-114mm)
    {不回流为\\开发反馈};
\end{tikzpicture}
